% !TEX root = ../../../main.tex
\section{Differentiation as a linear map}
\label{sec:linear-algebra-i:differentiation}

% Sources: supplied September 17 handout, slide 21, and Homework 4,
% Problem 5. The September 24 handout contains no differentiation topic.
% The author requested this Chapter 4 placement on October 1, 2026.
The elements of a vector space can be functions. The space \(C^0(\R)\)
consists of continuous real-valued functions on \(\R\), while \(C^1(\R)\)
consists of differentiable real-valued functions \(f\) for which both
\(f\) and \(f'\) are continuous. Both are vector spaces over \(\R\).
In this example, no finite-dimensional assumption is imposed.

Differentiation gives a map between these spaces,
\[
  D:C^1(\R)\longrightarrow C^0(\R),\qquad D(f)=f'.
\]
It is linear by the sum and constant factor rules of differentiation.
For \(f_1,f_2,f\in C^1(\R)\) and \(r\in\R\),
\[
  \begin{aligned}
    D(f_1+f_2)&=(f_1+f_2)'=f_1'+f_2'=D(f_1)+D(f_2),\\
    D(rf)&=(rf)'=rf'=rD(f).
  \end{aligned}
\]
Thus differentiation preserves addition and scalar multiplication.
Its inputs and outputs are functions.
Section~\ref{sec:linear-algebra-i:linear-combinations-and-bases}
discusses bases beyond finite dimension, including function spaces.

% Keep the complete shared polynomial exercise together in the book.
\newpage
\subsection{Polynomial differentiation}

The following exercise applies differentiation to finite-dimensional
polynomial spaces. It brings together the kernel, image, rank, nullity,
and matrix construction in chosen bases.

% Source: MA-GY7033 Fall 2026 Homework 4.pdf, p. 2, Problem 5.
\begin{exercise}
  \label{ex:hw04:polynomial-differentiation}
  \solutionlink{hw04:polynomial-differentiation}
  Let \(\mathcal P_3(\R)\) be the vector space of real polynomials of
  degree at most \(3\), including the zero polynomial, and define
  \(\mathcal P_2(\R)\) similarly. Consider the differentiation map
  \[
    D:\mathcal P_3(\R)\to\mathcal P_2(\R),\qquad D(p)=p'.
  \]
  \begin{enumerate}[(a)][2]
    \item Find \(\ker D\), \(\operatorname{im}D\), the nullity of \(D\),
      and the rank of \(D\).
    \item Determine whether \(D\) is injective and whether it is surjective.
    \item Find the matrix of \(D\) with respect to
      \[
        E=\begin{bmatrix}1&x&x^2&x^3\end{bmatrix},
        \qquad F=\begin{bmatrix}1&x&x^2\end{bmatrix}.
      \]
    \item Find a new ordered basis \(E'\) of \(\mathcal P_3(\R)\),
      beginning with a basis of \(\ker D\), for which the matrix relative
      to \(E'\) and \(F\) is
      \[
        \begin{bmatrix}
          0&1&0&0\\
          0&0&1&0\\
          0&0&0&1
        \end{bmatrix}.
      \]
  \end{enumerate}
\end{exercise}

